% GENERATED FILE. Edit canonical lessons, then run npm run generate:books.
% canonical-source-hash: fnv1a64:404f49af30b7a5ea
% canonical-lessons: HI-C541-rukavat, HI-C541-jhamela, HI-C541-atkana, HI-C541-phanda, HI-C541-uljhana

\chapter{Something in the Way}
\label{ch:hi-something-in-the-way}
\hlchaptermodality{541}

\begin{chapteropening}
\textbf{By the end of this chapter:} \emph{I can say what got in the way and who held things up.}

Complete the last lesson of chapter 541: I can say what got in the way and who held things up.
\end{chapteropening}

\section[\texorpdfstring{\hi{रुकावट} (rukāvaṭ)}{रुकावट (rukāvaṭ)}]{\hi{रुकावट} (rukāvaṭ) — an obstacle, a hold-up}

\label{lesson:HI-C541-rukavat}



\begin{quote}
Before the new one: say the Hindi for a trick, then the Hindi for an effort, an endeavour (formal).
\end{quote}

\subsection*{You'll want to know: \hi{रुकावट}}
\textbf{\hi{रुकावट}} — \emph{rukāvaṭ} — \textquotedblleft{}an obstacle, a hold-up\textquotedblright{}.

\begin{quote}
\hi{सड़क} \hi{पर} \hi{एक} \hi{पेड़} \hi{गिरा} \hi{था}, \hi{इसलिए} \hi{रास्ते} \hi{में} \hi{रुकावट} \hi{थी।}
\end{quote}

\begin{quote}
\emph{Saṛak par ek peṛ girā thā, isliye rāste meṁ rukāvaṭ thī.}
\end{quote}

\begin{quote}
A tree had fallen on the road, so the way was blocked.
\end{quote}

Feminine.

\begin{cousinweb}
\textbf{\hi{रुकावट}} (\emph{rukāvaṭ}) is the noun of \emph{ruknā}, \textquotedblleft{}to stop\textquotedblright{}, which CDIAL lists beside \emph{roknā} under a reconstructed \emph{*rokk} (CDIAL 10827), with the ending -āvaṭ that makes nouns of verbs.
\end{cousinweb}

\begin{grammarlens}[title={What you've built}]
One more word for saying what happens if.
\end{grammarlens}

\subsection*{Your turn}
\begin{itemize}
\raggedright
  \item \emph{Say it:} \emph{rukāvaṭ}
  \item \emph{Say it:} \emph{rukāvaṭ}, once more
  \item \emph{Say it:} \emph{Saṛak par ek peṛ girā thā, isliye rāste meṁ rukāvaṭ thī.} — A tree had fallen on the road, so the way was blocked.
  \item \emph{Recall:} say the Hindi for a trick, then the Hindi for an effort, an endeavour (formal), then say \emph{rukāvaṭ} again
\end{itemize}

\begin{tcolorbox}[breakable,colback=teal!4,colframe=teal!35!black,title={Before you move on}]
What does \hi{रुकावट} mean? (\textquotedblleft{}An obstacle, a hold-up\textquotedblright{}.) Say it once more.
\end{tcolorbox}
\section[\texorpdfstring{\hi{झमेला} (jhamelā)}{झमेला (jhamelā)}]{\hi{झमेला} (jhamelā) — a hassle, a tiresome business}

\label{lesson:HI-C541-jhamela}



\begin{quote}
Before the new one: say the Hindi for an effort, an endeavour (formal), then the Hindi for an obstacle, a hold-up.
\end{quote}

\subsection*{You'll want to know: \hi{झमेला}}
\textbf{\hi{झमेला}} — \emph{jhamelā} — \textquotedblleft{}a hassle, a tiresome business\textquotedblright{}.

\begin{quote}
\hi{दफ़्तर} \hi{में} \hi{काग़ज़ों} \hi{का} \hi{बड़ा} \hi{झमेला} \hi{है।}
\end{quote}

\begin{quote}
\emph{Daftar meṁ kāġzoṁ kā baṛā jhamelā hai.}
\end{quote}

\begin{quote}
There is a big hassle with papers at the office.
\end{quote}

Masculine.

\begin{cousinweb}
\textbf{\hi{झमेला}} (\emph{jhamelā}) is a tadbhava, inherited from \emph{*jhamm}, \textquotedblleft{}entanglement\textquotedblright{}, a form CDIAL reconstructs from its descendants (CDIAL 5340). CDIAL gives the Hindi word as \textquotedblleft{}wrangling\textquotedblright{}.
\end{cousinweb}

\begin{grammarlens}[title={What you've built}]
One more word for saying what happens if.
\end{grammarlens}

\subsection*{Your turn}
\begin{itemize}
\raggedright
  \item \emph{Say it:} \emph{jhamelā}
  \item \emph{Say it:} \emph{jhamelā}, once more
  \item \emph{Say it:} \emph{Daftar meṁ kāġzoṁ kā baṛā jhamelā hai.} — There is a big hassle with papers at the office.
  \item \emph{Recall:} say the Hindi for an effort, an endeavour (formal), then the Hindi for an obstacle, a hold-up, then say \emph{jhamelā} again
\end{itemize}

\begin{tcolorbox}[breakable,colback=teal!4,colframe=teal!35!black,title={Before you move on}]
What does \hi{झमेला} mean? (\textquotedblleft{}A hassle, a tiresome business\textquotedblright{}.) Say it once more.
\end{tcolorbox}
\section[\texorpdfstring{\hi{अटकाना} (aṭkānā)}{अटकाना (aṭkānā)}]{\hi{अटकाना} (aṭkānā) — to get (something) stuck; to hold up}

\label{lesson:HI-C541-atkana}



\begin{quote}
Before the new one: say the Hindi for an obstacle, a hold-up, then the Hindi for a hassle, a tiresome business.
\end{quote}

\subsection*{You'll want to know: \hi{अटकाना}}
\textbf{\hi{अटकाना}} — \emph{aṭkānā} — \textquotedblleft{}to get (something) stuck; to hold up\textquotedblright{}.

\begin{quote}
\hi{उसने} \hi{पतंग} \hi{पेड़} \hi{पर} \hi{अटका} \hi{दी।}
\end{quote}

\begin{quote}
\emph{Usne patang peṛ par aṭkā dī.}
\end{quote}

\begin{quote}
He got the kite stuck in the tree.
\end{quote}

Past, with \hi{ने} (he or she …): \hi{उसने} \hi{काम} \hi{अटकाया}, \emph{usne kām aṭkāyā}.

\begin{cousinweb}
\textbf{\hi{अटकाना}} (\emph{aṭkānā}) is the causative of \emph{aṭaknā}, \textquotedblleft{}to stick, to stop\textquotedblright{}, which CDIAL files under \emph{*aṭṭakk}, a reconstructed form (CDIAL 182).
\end{cousinweb}

\begin{grammarlens}[title={What you've built}]
One more word for saying what happens if.
\end{grammarlens}

\subsection*{Your turn}
\begin{itemize}
\raggedright
  \item \emph{Say it:} \emph{aṭkānā}
  \item \emph{Say it:} \emph{aṭkānā}, once more
  \item \emph{Say it:} \emph{Usne patang peṛ par aṭkā dī.} — He got the kite stuck in the tree.
  \item \emph{Recall:} say the Hindi for an obstacle, a hold-up, then the Hindi for a hassle, a tiresome business, then say \emph{aṭkānā} again
\end{itemize}

\begin{tcolorbox}[breakable,colback=teal!4,colframe=teal!35!black,title={Before you move on}]
What does \hi{अटकाना} mean? (\textquotedblleft{}To get (something) stuck; to hold up\textquotedblright{}.) Say it once more.
\end{tcolorbox}
\section[\texorpdfstring{\hi{फंदा} (phandā)}{फंदा (phandā)}]{\hi{फंदा} (phandā) — a noose; a snare}

\label{lesson:HI-C541-phanda}



\begin{quote}
Before the new one: say the Hindi for a hassle, a tiresome business, then the Hindi for to get (something) stuck.
\end{quote}

\subsection*{You'll want to know: \hi{फंदा}}
\textbf{\hi{फंदा}} — \emph{phandā} — \textquotedblleft{}a noose; a snare\textquotedblright{}.

\begin{quote}
\hi{किसान} \hi{ने} \hi{चूहों} \hi{के} \hi{लिए} \hi{एक} \hi{फंदा} \hi{रखा।}
\end{quote}

\begin{quote}
\emph{Kisān ne cūhoṁ ke liye ek phandā rakhā.}
\end{quote}

\begin{quote}
The farmer set a snare for the rats.
\end{quote}

Masculine.

\begin{cousinweb}
\textbf{\hi{फंदा}} (\emph{phandā}) is probably related to Hindi \emph{phā̃d} and \emph{phā̃s}, \textquotedblleft{}a noose, a net\textquotedblright{}, which CDIAL files under Sanskrit \emph{spāśa}, \textquotedblleft{}a fetter, a noose\textquotedblright{} (CDIAL 13813); CDIAL does not list \emph{phandā} itself.
\end{cousinweb}

\begin{grammarlens}[title={What you've built}]
One more word for saying what happens if.
\end{grammarlens}

\subsection*{Your turn}
\begin{itemize}
\raggedright
  \item \emph{Say it:} \emph{phandā}
  \item \emph{Say it:} \emph{phandā}, once more
  \item \emph{Say it:} \emph{Kisān ne cūhoṁ ke liye ek phandā rakhā.} — The farmer set a snare for the rats.
  \item \emph{Recall:} say the Hindi for a hassle, a tiresome business, then the Hindi for to get (something) stuck, then say \emph{phandā} again
\end{itemize}

\begin{tcolorbox}[breakable,colback=teal!4,colframe=teal!35!black,title={Before you move on}]
What does \hi{फंदा} mean? (\textquotedblleft{}A noose; a snare\textquotedblright{}.) Say it once more.
\end{tcolorbox}
\section[\texorpdfstring{\hi{उलझाना} (uljhānā)}{उलझाना (uljhānā)}]{\hi{उलझाना} (uljhānā) — to tangle; to confuse}

\label{lesson:HI-C541-uljhana}



\begin{quote}
Before the new one: say the Hindi for to get (something) stuck, then the Hindi for a noose.
\end{quote}

\subsection*{You'll want to know: \hi{उलझाना}}
\textbf{\hi{उलझाना}} — \emph{uljhānā} — \textquotedblleft{}to tangle; to confuse\textquotedblright{}.

\begin{quote}
\hi{बच्चे} \hi{ने} \hi{धागा} \hi{उलझा} \hi{दिया।}
\end{quote}

\begin{quote}
\emph{Bacce ne dhāgā uljhā diyā.}
\end{quote}

\begin{quote}
The child tangled up the thread.
\end{quote}

Past, with \hi{ने} (he or she …): \hi{उसने} \hi{बात} \hi{उलझाई}, \emph{usne bāt uljhāī}.

\begin{cousinweb}
\textbf{\hi{उलझाना}} (\emph{uljhānā}) is the causative of \emph{ulajhnā}, \textquotedblleft{}to get tangled\textquotedblright{}, which CDIAL files under Sanskrit \emph{uparundhati}, \textquotedblleft{}obstructs\textquotedblright{} (CDIAL 2221).
\end{cousinweb}

\begin{grammarlens}[title={What you've built}]
One more word for saying what happens if.
\end{grammarlens}

\subsection*{Your turn}
\begin{itemize}
\raggedright
  \item \emph{Say it:} \emph{uljhānā}
  \item \emph{Say it:} \emph{uljhānā}, once more
  \item \emph{Say it:} \emph{Bacce ne dhāgā uljhā diyā.} — The child tangled up the thread.
  \item \emph{Recall:} say the Hindi for to get (something) stuck, then the Hindi for a noose, then say \emph{uljhānā} again
\end{itemize}

\begin{tcolorbox}[breakable,colback=teal!4,colframe=teal!35!black,title={Before you move on}]
What does \hi{उलझाना} mean? (\textquotedblleft{}To tangle; to confuse\textquotedblright{}.) Say it once more.
\end{tcolorbox}
